% =====================================================================
%  Chapitre 9 - Sprint 6 : collaboration et aide a la decision
% =====================================================================
\chapter{Sprint 6 : collaboration et aide a la decision}
\label{chap:sprint6}

\section*{Introduction}
\addcontentsline{toc}{section}{Introduction}

Ce dernier sprint traite de trois services qui ne relevent d'aucun
domaine metier au sens strict, mais qui conditionnent l'usage quotidien
de la plateforme : la messagerie interne, l'assistant documentaire et le
moteur d'optimisation de flotte. Les deux derniers completent l'ensemble
des quatre services d'aide a la decision, dont l'estimation tarifaire et
la prediction de pannes ont deja ete traitees.

% =====================================================================
\section{Planification du sprint}

\subsection{Objectifs}

\begin{enumerate}
  \item Doter la plateforme d'une messagerie interne en temps reel.
  \item Offrir au client un assistant capable de repondre dans sa langue.
  \item Assister le gestionnaire dans l'affectation des vehicules et des
    chauffeurs.
  \item Consolider les indicateurs d'exploitation dans un tableau de
    bord.
\end{enumerate}

\subsection{Carnet du sprint}

\begin{table}[H]
\centering
\caption{Carnet du sprint 6}
\label{tab:backlog-s6}
\small
\begin{tabularx}{\textwidth}{@{}cp{2.3cm}Xcc@{}}
\toprule
\textbf{N\textsuperscript{o}} & \textbf{Acteur} & \textbf{Besoin exprime} &
\textbf{Prio.} & \textbf{Pts} \\ \midrule
1 & Utilisateur & Echanger en conversation privee ou en groupe & Should & 8 \\
2 & Utilisateur & Joindre un fichier a un message & Should & 5 \\
3 & Utilisateur & Voir qui est connecte & Could & 3 \\
4 & Client & Interroger un assistant sur les services & Could & 8 \\
5 & Client & Obtenir une reponse dans sa propre langue & Could & 5 \\
6 & Administrateur & Reindexer la base documentaire & Could & 3 \\
7 & Responsable & Obtenir une proposition d'affectation & Could & 8 \\
8 & Responsable & Consulter le tableau de bord consolide & Should & 5 \\
\midrule
\multicolumn{4}{@{}r}{\textbf{Total}} & \textbf{45} \\
\bottomrule
\end{tabularx}
\end{table}

% =====================================================================
\section{Messagerie interne}

\figReelle{chat-usecase.png}
  {Cas d'utilisation du service de messagerie}
  {fig:uc-chat}{1.00}

\figReelle{chat-sequence.png}
  {Sequence : envoi d'un message en temps reel}
  {fig:seq-chat}{1.00}

La messagerie est le seul service de la plateforme a maintenir une
\essentiel{liaison permanente} avec le navigateur. La
figure~\ref{fig:seq-chat} en montre la consequence : un message envoye
est d'abord persiste, puis diffuse aux participants connectes sans que
leur navigateur ait a interroger le serveur.

Ce choix a une contrepartie sur l'infrastructure. La passerelle doit
\essentiel{laisser passer la negociation de protocole} plutot que la
traiter comme une requete ordinaire, faute de quoi la liaison retombe sur
un mode degrade. C'est le seul endroit du systeme ou la configuration du
proxy porte une regle specifique.

\figReelle{chat-class-domain.png}
  {Modele du domaine du service de messagerie}
  {fig:dom-chat}{0.60}

Le modele repose sur trois classes. \exemple{Conversation} porte le type,
privee ou de groupe. \exemple{Participant} est une classe d'association
entre la conversation et l'utilisateur : elle porte la date de derniere
lecture, l'epinglage et l'activation des notifications, c'est-a-dire
\essentiel{ce qui est propre a chaque membre} et non a la conversation.
\exemple{Message} porte le contenu, un eventuel fichier joint et les
marqueurs de lecture.

% =====================================================================
\section{Assistant documentaire}

\figReelle{ayshaa-rag-usecase.png}
  {Cas d'utilisation de l'assistant documentaire}
  {fig:uc-rag}{1.00}

\figReelle{ayshaa-rag-sequence.png}
  {Sequence : traitement d'une question posee a l'assistant}
  {fig:seq-rag}{1.00}

L'assistant repond a partir d'une base documentaire vectorisee. La
sequence de la figure~\ref{fig:seq-rag} suit quatre etapes : detection de
la langue, recherche dans la base, consultation eventuelle des donnees
metier, puis redaction de la reponse.

\paragraph{Une recherche hybride.}
La recherche combine deux methodes. La premiere compare des
\essentiel{vecteurs de sens} : elle retrouve un document dont le
vocabulaire differe de celui de la question. La seconde est
\essentiel{lexicale} : elle privilegie la correspondance exacte des
termes. Les deux se completent, la premiere rattrapant les reformulations
et la seconde les references precises comme un numero de facture.

\paragraph{Trois langues, dont le dialecte local.}
Le detecteur de langue distingue le francais, l'anglais et l'arabe, en
prenant en charge \essentiel{le dialecte tunisien et sa transcription en
caracteres latins}. Ce point a justifie l'ecriture d'un detecteur propre
au projet : les detecteurs generiques traitent mal des messages courts et
melangeant les langues, ce qui est la norme dans les echanges
commerciaux locaux.

\paragraph{Un repli assume.}
La redaction finale s'appuie sur un modele de generation lorsqu'il est
disponible. \essentiel{S'il ne l'est pas, l'assistant ne tombe pas en
panne} : il restitue les extraits documentaires les plus pertinents. La
reponse est moins fluide, mais elle reste exacte et sourcee.

% =====================================================================
\section{Optimisation de l'affectation}

\figReelle{ai-optimizer-usecase.png}
  {Cas d'utilisation du moteur d'optimisation}
  {fig:uc-optimizer}{1.00}

\figReelle{ai-optimizer-sequence.png}
  {Sequence : production d'un plan d'affectation}
  {fig:seq-optimizer}{1.00}

Le moteur d'optimisation propose une affectation de vehicules et de
chauffeurs aux missions en attente. Il procede en deux temps
\essentiel{qu'il importe de distinguer} : d'abord il ecarte les
candidatures impossibles, ensuite il note celles qui restent.

\begin{table}[H]
\centering
\caption{Contraintes eliminatoires et criteres de notation}
\label{tab:optimisation}
\small
\begin{tabularx}{\textwidth}{@{}p{3.6cm}X@{}}
\toprule
\textbf{Nature} & \textbf{Regles appliquees} \\ \midrule
Contraintes eliminatoires & Classe du vehicule adaptee au chargement,
  categorie de permis du chauffeur, vehicule non immobilise pour
  maintenance, absence de chevauchement d'horaires, amplitude de conduite
  respectee \\
Criteres de notation & Proximite, disponibilite, adequation du vehicule a
  la mission, equilibrage de la charge entre chauffeurs \\
\bottomrule
\end{tabularx}
\end{table}

\essentiel{Le moteur propose, il ne decide pas.} Le plan produit est
soumis au gestionnaire, qui l'applique ou le modifie. Ce parti pris est le
meme que celui retenu pour l'estimation tarifaire au sprint~3 : l'aide a
la decision reste une aide.

L'echange est \essentiel{asynchrone} : la demande d'optimisation est
publiee sur le courtier de messages, le moteur la consomme et republie le
plan. Le gestionnaire n'attend pas devant son ecran, ce qui est acceptable
pour une operation de planification et impossible pour une estimation de
prix.

% =====================================================================
\section{Realisation}

\subsection{Les outils de collaboration}

La messagerie interne, figure~\ref{fig:s6-messagerie}, remplace les
groupes de discussion externes que l'entreprise utilisait pour coordonner
ses equipes. Le gain n'est pas le confort mais le \essentiel{rattachement
au contexte metier} : une conversation est liee a une mission, a un
vehicule ou a un ordre de travail, si bien que l'echange reste retrouvable
depuis l'objet qu'il concerne. Sur un fil de discussion generaliste, la
meme information se perd en quelques jours sous le flux des messages
suivants.

\capture{s6-messagerie.jpg}
  {Messagerie interne rattachee au contexte metier}
  {fig:s6-messagerie}

L'assistant documentaire du portail client, figure~\ref{fig:s6-assistant},
repond aux questions courantes sur les services proposes, les zones
desservies et les modalites de reservation. Il s'appuie sur les documents
de l'entreprise et \essentiel{cite ses sources}, ce qui permet au visiteur
de verifier la reponse. Son interet pour l'exploitation est de filtrer les
sollicitations les plus repetitives, qui occupaient jusque-la une part
non negligeable du temps telephonique.

\capture{s6-assistant-ayshaa.jpg}
  {Assistant documentaire du portail client}
  {fig:s6-assistant}

\subsection{Les outils d'aide a la decision}

Le plan d'affectation, figure~\ref{fig:s6-optimisateur}, restitue le
travail du moteur decrit a la section precedente. Chaque mission en
attente y recoit une proposition de vehicule et de chauffeur, assortie
de la note obtenue et des criteres qui l'ont determinee. Le gestionnaire
conserve la main : il applique le plan tel quel, ou il substitue une
affectation. Cette reversibilite est ce qui rend l'outil acceptable en
exploitation, ou une contrainte connue du seul exploitant --- un client
difficile, un chauffeur en formation --- echappe par nature au moteur.

\capture{s6-optimisateur.jpg}
  {Plan d'affectation propose par le moteur d'optimisation}
  {fig:s6-optimisateur}

Le tableau de bord, figure~\ref{fig:s6-bord1}, consolide en une page les
indicateurs d'exploitation : missions en cours et terminees, taux
d'utilisation du parc, vehicules immobilises, alertes ouvertes. Il agrege
des donnees produites par plusieurs services, ce qui en fait le point de
la plateforme ou la separation en services independants est la moins
visible pour l'utilisateur --- et c'est bien le resultat recherche.

\capture{s6-tableau-bord-1.jpg}
  {Tableau de bord : indicateurs d'exploitation consolides}
  {fig:s6-bord1}

% =====================================================================
\section{Tests de validation}

\begin{table}[H]
\centering
\caption{Tests de validation du sprint 6}
\label{tab:tests-s6}
\small
\begin{tabularx}{\textwidth}{@{}cXp{2.9cm}c@{}}
\toprule
\textbf{N\textsuperscript{o}} & \textbf{Cas teste} & \textbf{Attendu} &
\textbf{Etat} \\ \midrule
T1 & Message envoye a un groupe & Recu sans rechargement de page & Reussi \\
T2 & Fichier joint a un message & Telechargeable par les participants & Reussi \\
T3 & Question posee en dialecte tunisien & Reponse dans la meme langue & Reussi \\
T4 & Modele de generation indisponible & Extraits documentaires restitues & Reussi \\
T5 & Chauffeur sans le permis requis & Ecarte du plan d'affectation & Reussi \\
T6 & Vehicule en maintenance & Ecarte du plan d'affectation & Reussi \\
\bottomrule
\end{tabularx}
\end{table}

% =====================================================================
\section*{Conclusion}
\addcontentsline{toc}{section}{Conclusion}

Ce dernier sprint livre la messagerie interne, l'assistant documentaire
et le moteur d'optimisation. Les quarante-cinq points du carnet ont ete
realises, ce qui porte a trois cent six le total des points du projet.

Une ligne de conduite commune relie les quatre services d'aide a la
decision developpes au fil des sprints : \essentiel{aucun ne decide a la
place de l'utilisateur}. L'estimation propose un prix que le gestionnaire
peut modifier, la prediction signale un risque en le justifiant, l'assistant
cite ses sources, l'optimisation soumet un plan. Cette retenue n'est pas
une limite technique mais un choix : dans une entreprise de transport, la
responsabilite d'immobiliser un vehicule ou d'engager un tarif reste
humaine.

La realisation etant achevee, la conclusion generale dresse le bilan du
projet et en ouvre les perspectives.
